% GENERATED FILE. Edit canonical lessons, then run npm run generate:books.
% canonical-source-hash: fnv1a64:b50c21acc0741026
% canonical-lessons: RU-C83-midday, RU-C83-midnight, RU-C83-early, RU-C83-late, RU-C83-in-the-morning

\chapter{Midday, Midnight, Early, Late, In the morning}
\label{ch:ru-midday-midnight-early-late-in-the-morning}
\hlchaptermodality{83}

\begin{chapteropening}
\textbf{By the end of this chapter:} \emph{I can say midday, midnight, early, late and in the morning.}

Complete the last lesson of chapter 83: I can say midday, midnight, early, late and in the morning.
\end{chapteropening}

\section[\texorpdfstring{\ru{полдень} (pólden)}{полдень (pólden)}]{\ru{полдень} (pólden) — midday}

\label{lesson:RU-C83-midday}



\begin{quote}
Before the new one: say the Russian for that, then the Russian for such.
\end{quote}

\subsection*{You'll want to know: \ru{полдень}}
\textbf{\ru{полдень}} (\emph{pólden}) — \textquotedblleft{}midday\textquotedblright{}.

\begin{cousinweb}
\textbf{\ru{полдень}} is \emph{\ru{пол}-}, \textquotedblleft{}half\textquotedblright{}, with \emph{\ru{день}}, \textquotedblleft{}day\textquotedblright{}: the middle of the day.
\end{cousinweb}

\begin{grammarlens}[title={What you've built}]
One more word for this chapter.
\end{grammarlens}

\subsection*{Your turn}
\begin{itemize}
\raggedright
  \item \emph{Say it:} \emph{pólden}
  \item \emph{Say it:} \emph{pólden}, once more
  \item \emph{Say it:} say \emph{pólden}
  \item \emph{Recall:} say the Russian for that, then the Russian for such, then say \emph{pólden} again
\end{itemize}

\begin{tcolorbox}[breakable,colback=teal!4,colframe=teal!35!black,title={Before you move on}]
What does \ru{полдень} mean? (\textquotedblleft{}Midday\textquotedblright{}.) Say it once more.
\end{tcolorbox}
\section[\texorpdfstring{\ru{полночь} (pólnoch)}{полночь (pólnoch)}]{\ru{полночь} (pólnoch) — midnight}

\label{lesson:RU-C83-midnight}



\begin{quote}
Before the new one: say the Russian for such, then the Russian for midday.
\end{quote}

\subsection*{You'll want to know: \ru{полночь}}
\textbf{\ru{полночь}} (\emph{pólnoch}) — \textquotedblleft{}midnight\textquotedblright{}.

\begin{cousinweb}
\textbf{\ru{полночь}} is \emph{\ru{пол}-}, \textquotedblleft{}half\textquotedblright{}, with \emph{\ru{ночь}}, \textquotedblleft{}night\textquotedblright{}. \emph{\ru{ночь}} $\leftarrow$ Proto-Slavic \emph{*noťĭ}, a cousin of Lithuanian \emph{naktìs}, Latin \emph{nox} and English \emph{night} (Vasmer).
\end{cousinweb}

\begin{grammarlens}[title={What you've built}]
One more word for this chapter.
\end{grammarlens}

\subsection*{Your turn}
\begin{itemize}
\raggedright
  \item \emph{Say it:} \emph{pólnoch}
  \item \emph{Say it:} \emph{pólnoch}, once more
  \item \emph{Say it:} say \emph{pólnoch}
  \item \emph{Recall:} say the Russian for such, then the Russian for midday, then say \emph{pólnoch} again
\end{itemize}

\begin{tcolorbox}[breakable,colback=teal!4,colframe=teal!35!black,title={Before you move on}]
What does \ru{полночь} mean? (\textquotedblleft{}Midnight\textquotedblright{}.) Say it once more.
\end{tcolorbox}
\section[\texorpdfstring{\ru{рано} (ráno)}{рано (ráno)}]{\ru{рано} (ráno) — early}

\label{lesson:RU-C83-early}



\begin{quote}
Before the new one: say the Russian for midday, then the Russian for midnight.
\end{quote}

\subsection*{You'll want to know: \ru{рано}}
\textbf{\ru{рано}} (\emph{ráno}) — \textquotedblleft{}early\textquotedblright{}.

\begin{cousinweb}
\textbf{\ru{рано}} $\leftarrow$ Proto-Slavic \emph{*rano}, \textquotedblleft{}early\textquotedblright{} — inherited, but Vasmer calls it hard to explain: some link it with Lithuanian \emph{rýtas}, \textquotedblleft{}morning\textquotedblright{}, others with Greek \emph{órthros}, \textquotedblleft{}dawn\textquotedblright{}.
\end{cousinweb}

\begin{grammarlens}[title={What you've built}]
One more word for this chapter.
\end{grammarlens}

\subsection*{Your turn}
\begin{itemize}
\raggedright
  \item \emph{Say it:} \emph{ráno}
  \item \emph{Say it:} \emph{ráno}, once more
  \item \emph{Say it:} say \emph{ráno}
  \item \emph{Recall:} say the Russian for midday, then the Russian for midnight, then say \emph{ráno} again
\end{itemize}

\begin{tcolorbox}[breakable,colback=teal!4,colframe=teal!35!black,title={Before you move on}]
What does \ru{рано} mean? (\textquotedblleft{}Early\textquotedblright{}.) Say it once more.
\end{tcolorbox}
\section[\texorpdfstring{\ru{поздно} (pózdno)}{поздно (pózdno)}]{\ru{поздно} (pózdno) — late}

\label{lesson:RU-C83-late}



\begin{quote}
Before the new one: say the Russian for midnight, then the Russian for early.
\end{quote}

\subsection*{You'll want to know: \ru{поздно}}
\textbf{\ru{поздно}} (\emph{pózdno}) — \textquotedblleft{}late\textquotedblright{}.

\begin{cousinweb}
\textbf{\ru{поздно}} is the adverb of \emph{\ru{поздний}}, \textquotedblleft{}late\textquotedblright{}, Proto-Slavic \emph{*pozdŭ} — inherited, from the PIE \emph{*pos}, \textquotedblleft{}after, behind\textquotedblright{}, which Vasmer also finds in Lithuanian \emph{pàs}, \textquotedblleft{}at\textquotedblright{}, and Latin \emph{post}, \textquotedblleft{}after\textquotedblright{}.
\end{cousinweb}

\begin{grammarlens}[title={What you've built}]
One more word for this chapter.
\end{grammarlens}

\subsection*{Your turn}
\begin{itemize}
\raggedright
  \item \emph{Say it:} \emph{pózdno}
  \item \emph{Say it:} \emph{pózdno}, once more
  \item \emph{Say it:} say \emph{pózdno}
  \item \emph{Recall:} say the Russian for midnight, then the Russian for early, then say \emph{pózdno} again
\end{itemize}

\begin{tcolorbox}[breakable,colback=teal!4,colframe=teal!35!black,title={Before you move on}]
What does \ru{поздно} mean? (\textquotedblleft{}Late\textquotedblright{}.) Say it once more.
\end{tcolorbox}
\section[\texorpdfstring{\ru{утром} (útrom)}{утром (útrom)}]{\ru{утром} (útrom) — in the morning}

\label{lesson:RU-C83-in-the-morning}



\begin{quote}
Before the new one: say the Russian for early, then the Russian for late.
\end{quote}

\subsection*{You'll want to know: \ru{утром}}
\textbf{\ru{утром}} (\emph{útrom}) — \textquotedblleft{}in the morning\textquotedblright{}.

\begin{cousinweb}
\textbf{\ru{утром}} is \emph{\ru{утро}}, \textquotedblleft{}morning\textquotedblright{}, in the form that says \textquotedblleft{}by, during\textquotedblright{}: \textquotedblleft{}in the morning\textquotedblright{}. \emph{\ru{утро}} $\leftarrow$ Proto-Slavic \emph{*jutro}; Vasmer calls its origin difficult, with links to \emph{\ru{уже}}, \textquotedblleft{}already\textquotedblright{}, or to Latin \emph{iugum}, \textquotedblleft{}yoke\textquotedblright{}, as \textquotedblleft{}the time of yoking\textquotedblright{}.
\end{cousinweb}

\begin{grammarlens}[title={What you've built}]
One more word for this chapter.
\end{grammarlens}

\subsection*{Your turn}
\begin{itemize}
\raggedright
  \item \emph{Say it:} \emph{útrom}
  \item \emph{Say it:} \emph{útrom}, once more
  \item \emph{Say it:} say \emph{útrom}
  \item \emph{Recall:} say the Russian for early, then the Russian for late, then say \emph{útrom} again
\end{itemize}

\begin{tcolorbox}[breakable,colback=teal!4,colframe=teal!35!black,title={Before you move on}]
What does \ru{утром} mean? (\textquotedblleft{}In the morning\textquotedblright{}.) Say it once more.
\end{tcolorbox}
